\chapter{Relativité restreinte~: dilatation du temps}\label{ch:g12:special-relativity}

Chaque \omterm{def:g10:universal-gravitation:satellite}{satellite} GPS porte une horloge atomique bonne à une
seconde en trois millions d'années --- et avant le lancement, les
ingénieurs la désaccordent exprès~: laissée honnête, elle gagnerait
$38$~millionièmes de seconde par jour, et demain soir votre téléphone
vous placerait à onze kilomètres de là où vous êtes. L'horloge n'est
pas en faute --- le temps lui-même tic-tac différemment pour le mobile
et l'immobile~: ce chapitre suit Einstein d'un fait obstiné à cette
conclusion.

\section{Une vitesse refuse de s'ajouter}

Les vitesses dépendent du référentiel et se composent par addition
(\cref{ch:g10:relative-motion})~: un passager marchant à \qty{4}{km/h}
vers l'avant d'un train à \qty{300}{km/h} passe l'observateur au bord
de la voie à \qty{304}{km/h}.

\begin{definition}[Référentiel galiléen]\label{def:g12:special-relativity:inertial}
Un \emph{référentiel galiléen}\index{référentiel galiléen} est un
\omterm{def:g10:relative-motion:frame}{référentiel} (\cref{ch:g10:relative-motion}) dans lequel un
corps sans force garde une vitesse constante~: le sol, à peu près, et
tout référentiel en mouvement rectiligne uniforme par rapport à lui.
\end{definition}

Pousser la règle à l'extrême~: le faisceau de phare du train devrait
passer l'observateur au sol à $c + \qty{300}{km/h}$. En 1887 Michelson
et Morley ont comparé la vitesse de la lumière le long et en travers
de la course orbitale de \qty{30}{km/s} de la Terre, chassant de telles
différences~; à une précision exquise ils n'en ont trouvé aucune ---
le même $c$ dans chaque direction, en chaque saison. Einstein, en
1905, a cru l'expérience et reconstruit le temps et l'espace autour
d'elle, sur deux postulats jugés par leurs conséquences --- testés
depuis un siècle sans un seul échec.

\begin{theorem}[Les postulats de la relativité restreinte]\label{thm:g12:special-relativity:postulates}
\index{relativité restreinte}\index{postulats de la relativité}
\begin{enumerate}
  \item \emph{Principe de relativité}~: les lois de la physique
        prennent la même forme dans tout
        \omterm{def:g12:special-relativity:inertial}{référentiel galiléen}~; aucun n'est
        «~vraiment au repos~».
  \item \emph{Invariance de $c$}~: la lumière dans le vide voyage à
        $c = \qty{3.00e8}{m/s}$ dans tout
        \omterm{def:g12:special-relativity:inertial}{référentiel galiléen}, quel que soit le
        mouvement de la source ou de l'observateur.
\end{enumerate}
\end{theorem}
\admitted

\section{Adieu au «~même instant~»}

\begin{proposition}[Relativité de la simultanéité]\label{prop:g12:special-relativity:simultaneity}
\index{simultanéité}
Deux \omterm{def:g12:special-relativity:proper-time}{événements} simultanés dans un
\omterm{def:g12:special-relativity:inertial}{référentiel galiléen} ne le sont, en général, pas
dans un autre~: «~maintenant, partout~» dépend du référentiel.
\end{proposition}

\begin{proof}
Le train d'Einstein. Deux éclairs frappent les deux extrémités d'un
train en mouvement, scorchant la voie. L'observateur du quai $M$, à
mi-chemin entre les marques, reçoit les deux éclairs ensemble~;
chacun a traversé la même distance à la même vitesse $c$, donc les
frappes étaient simultanées. Le passager $M'$, à mi-chemin le long
du train, roule \emph{vers} l'éclair avant, qui l'atteint le
premier~; mais dans son référentiel aussi les deux éclairs ont
couvert des distances égales (la moitié d'un train) à la même
vitesse $c$ --- deuxième postulat~! --- donc la frappe avant s'est
produite la première. Chaque verdict est correct dans son propre
référentiel~; le désaccord, des nanosecondes pour un train, est
pourquoi personne ne l'avait remarqué.
\end{proof}

\begin{omfigure}
\begin{tikzpicture}[scale=0.82, font=\small]
  \draw[thick] (0,0) -- (10.6,0);
  \fill (5.3,0) circle (2pt); \node[below] at (5.3,-0.05) {$M$};
  \node[below] at (1.6,-0.05) {$A$}; \node[below] at (9.0,-0.05) {$B$};
  \draw[omDef, thick] (1.6,0.55) rectangle (9.0,1.45);
  \fill[omDef] (5.3,1.45) circle (2pt); \node[above] at (5.3,1.5) {$M'$};
  \draw[-Stealth, omDef, thick] (9.2,1.0) -- (10.4,1.0) node[above left] {$\vect v$};
  \draw[omThm, thick] (1.25,2.7) -- (1.7,2.15) -- (1.45,2.05) -- (1.75,1.5);
  \draw[omThm, thick] (8.65,2.7) -- (9.1,2.15) -- (8.85,2.05) -- (9.15,1.5);
  \draw[dashed, omExo] (2.35,0) arc (0:60:0.75) (3.1,0) arc (0:60:1.5);
  \draw[dashed, omExo] (8.25,0) arc (180:120:0.75) (7.5,0) arc (180:120:1.5);
\end{tikzpicture}

{\small Le train d'Einstein~: les éclairs de $A$ et $B$ atteignent $M$
ensemble~; $M'$ roule vers l'éclair de $B$, le rencontre le premier
--- et conclut qu'il a frappé le premier.}
\end{omfigure}

\section{L'horloge à lumière et le facteur $\gamma$}

\begin{definition}[Événement, temps propre]\label{def:g12:special-relativity:proper-time}
Un \emph{événement}\index{événement} est une occurrence en un lieu et
un instant --- un tic, un éclair, une désintégration. Lorsqu'une seule
horloge est présente à deux événements, l'intervalle qu'elle lit entre
eux est le \emph{temps propre}\index{temps propre} $\Delta t_0$. Un
référentiel dans lequel cette horloge se déplace mesure un intervalle
$\Delta t$ en général différent.
\end{definition}

\begin{theorem}[Dilatation du temps]\label{thm:g12:special-relativity:dilation}
\index{dilatation du temps}\index{facteur gamma@$\gamma$ (gamma) facteur}
Une horloge se déplaçant à vitesse constante $v$ à travers un
\omterm{def:g12:special-relativity:inertial}{référentiel galiléen} y est mesurée comme
retardant~: entre deux \omterm{def:g12:special-relativity:proper-time}{événements} à
l'horloge,
\[
\Delta t = \gamma\, \Delta t_0, \qquad
\gamma = \frac{1}{\sqrt{1 - v^2/c^2}} \geq 1 .
\]
\end{theorem}

\begin{proof}
Construire l'horloge la plus simple imaginable~: deux miroirs en
regard à une distance $L$, un \omterm{def:g12:quantum-world:photon}{photon} rebondissant entre
eux~; un aller-retour est un tic, et au repos un tic dure
$\Delta t_0 = 2L/c$. Maintenant boulonner l'horloge à un vaisseau
traversant à la vitesse $v$, miroirs perpendiculaires au mouvement, et
chronométrer un tic depuis le sol~: pendant que le
\omterm{def:g12:quantum-world:photon}{photon} grimpe, les miroirs glissent de côté, donc le
\omterm{def:g12:quantum-world:photon}{photon} vole un chemin oblique --- deux mouvements à la
fois, comme le \omterm{def:g12:projectile-motion:projectile}{projectile} du
\cref{ch:g12:kinematics-2d}. Si le tic dure $\Delta t$, chaque
demi-tic avance de $v\,\Delta t/2$ horizontalement et de $L$
\omterm{def:g10:universal-gravitation:weight}{verticalement}~; et --- l'étape cruciale --- le
deuxième postulat fixe la vitesse au sol du
\omterm{def:g12:quantum-world:photon}{photon} à $c$, donc le demi-chemin oblique mesure
$c\,\Delta t/2$. Pythagore~:
\[
\Bigl(\frac{c\,\Delta t}{2}\Bigr)^{\!2} = L^2 + \Bigl(\frac{v\,\Delta t}{2}\Bigr)^{\!2}
\;\Longrightarrow\;
\Delta t = \frac{2L/c}{\sqrt{1 - v^2/c^2}} = \gamma\,\Delta t_0 .
\]
Et toute horloge qui voyage à côté --- montre, battement de cœur,
\omterm{def:g11:fundamental-interactions:ladder}{noyau} qui se désintègre --- doit rester en
pas, ou la comparaison trahirait le mouvement du vaisseau et
violerait le premier postulat.
\end{proof}

\begin{omfigure}
\begin{tikzpicture}[scale=0.9, font=\small]
  \draw[thick] (-0.5,0) -- (0.5,0); \draw[thick] (-0.5,2) -- (0.5,2);
  \draw[omDef, thick, -Stealth] (0,0.05) -- (0,1.95);
  \draw[Stealth-Stealth] (0.85,0) -- (0.85,2) node[midway, right] {$L$};
  \node at (0.3,-0.75) {au repos~: tic $\Delta t_0 = 2L/c$};
\end{tikzpicture}
\qquad
\begin{tikzpicture}[scale=0.9, font=\small]
  \foreach \x in {0,1.4,2.8} { \draw[gray] (\x-0.4,0) -- (\x+0.4,0);
    \draw[gray] (\x-0.4,2) -- (\x+0.4,2); }
  \draw[omDef, thick, -Stealth] (0,0.03) -- (1.4,1.97);
  \draw[omDef, thick, -Stealth] (1.4,1.97) -- (2.8,0.03);
  \draw[dashed, omExo] (1.4,0) -- (1.4,2);
  \draw[-Stealth, thick] (3.5,1.0) -- (4.4,1.0) node[above left] {$\vect v$};
  \node[above left] at (0.9,0.9) {$\frac{c\,\Delta t}{2}$};
  \node[right] at (1.4,0.75) {$L$};
  \draw[Stealth-Stealth] (0,-0.35) -- (1.4,-0.35) node[midway, below] {$v\,\Delta t/2$};
  \node at (1.5,-1.15) {en mouvement~: le photon vole l'hypoténuse};
\end{tikzpicture}

{\small L'horloge à lumière. Au repos le
\omterm{def:g12:quantum-world:photon}{photon} grimpe $L$~; depuis le sol il vole
l'hypoténuse au même $c$ --- donc le tic doit prendre plus longtemps.}
\end{omfigure}

\begin{example}[Gamma en nombres]\label{ex:g12:special-relativity:gammatable}
\begin{center}\small
\begin{tabular}{lccccccc}
$v/c$    & 0.01     & 0.1   & 0.5   & 0.8   & 0.9   & 0.99 & 0.999\\
$\gamma$ & 1.00005 & 1.005 & 1.155 & 1.667 & 2.294 & 7.09 & 22.4
\end{tabular}
\end{center}
Rien ne se passe pendant des pages, puis tout~: $\gamma$ colle à $1$
jusqu'à un tiers de la vitesse de la lumière, passe $2$ près de
$0.9c$, et explose lorsque $v \to c$.
\end{example}

\begin{omfigure}
\begin{tikzpicture}
\begin{axis}[omaxis, width=8.6cm, height=5.4cm,
    xlabel={$v/c$}, ylabel={$\gamma$},
    xmin=0, xmax=1.08, ymin=0, ymax=8,
    xtick={0,0.2,0.4,0.6,0.8,1}, ytick={1,2,3,4,5,6,7}]
  \addplot[omDef, thick, domain=0:0.992, samples=200] {1/sqrt(1-x^2)};
  \addplot[omExo, dashed, thick] coordinates {(1,0) (1,8)};
\end{axis}
\end{tikzpicture}

{\small $\gamma = 1/\sqrt{1 - v^2/c^2}$~: indiscernable de $1$ aux
vitesses quotidiennes, \omterm{def:g11:lenses-and-eye:lens}{divergent} à l'asymptote $v = c$
--- la vitesse qu'aucun corps massif n'atteint.}
\end{omfigure}

\begin{remark}[Pourquoi la vie quotidienne n'a jamais remarqué]\label{rem:g12:special-relativity:everyday}
Pour $v \ll c$, écrire $x = v^2/c^2$~: puisque $(1-x)(1+x) \approx 1$
et $(1 + x/2)^2 \approx 1 + x$, on obtient $\gamma \approx 1 +
v^2/(2c^2)$. Une voiture à \qty{130}{km/h} a $\gamma - 1 \approx
\num{7e-15}$ (une seconde perdue pour quatre millions d'années de
conduite)~; un jet gère $\num{3.5e-13}$, une seconde pour quatre-vingt-dix
mille ans en l'air. Nos horloges étaient simplement trop grossières.
\end{remark}

\section{Le verdict de l'expérience}

\begin{example}[Les muons atteignent le sol]\label{ex:g12:special-relativity:muons}
Les rayons cosmiques frappant la haute atmosphère créent des
\emph{muons} --- particules de courte vie de durée de vie propre
$\Delta t_0 = \qty{2.2}{\micro s}$ --- à environ \qty{15}{km} de
hauteur, se déplaçant à près de $c$. Sans relativité un muon parcourt
$c\,\Delta t_0 \approx \qty{660}{m}$ par durée de vie~: la descente
prend quelque $23$~durées de vie, et la fraction survivante
$\mathrm{e}^{-23} \approx \num{e-10}$ devrait rendre les muons au
niveau de la mer infiniment rares. Pourtant les détecteurs en comptent
environ un par centimètre carré par minute. La
\omterm{thm:g12:special-relativity:dilation}{dilatation du temps} le résout~: arrivant avec
$\gamma \approx 30$, les muons vivent $30 \times \qty{2.2}{\micro s}
= \qty{66}{\micro s}$ sur nos horloges ---
\omterm{def:g12:projectile-motion:range}{portée} $\approx \qty{20}{km}$, toute l'atmosphère.
\end{example}

\begin{omfigure}
\begin{tikzpicture}[scale=0.85, font=\small]
  \draw[thick] (0,0) -- (8.6,0) node[below left] {sol};
  \draw[dashed, omExo] (0,5) -- (8.6,5);
  \node[above] at (2.1,5.02) {muons nés à \qty{15}{km}};
  \draw[-Stealth, thick] (1.1,6.1) node[above right] {rayon cosmique} -- (1.55,5.05);
  \draw[omThm, very thick, -Stealth] (3.4,5) -- (3.4,4.78);
  \node[right, omThm] at (3.5,4.62)
    {sans dilatation~: $c\,\Delta t_0 \approx \qty{660}{m}$};
  \draw[omDef, very thick, -Stealth] (6.6,5) -- (6.6,0.06);
  \node[left, omDef, align=right] at (6.5,2.2)
    {$\gamma \approx 30$~:\\ $\approx \qty{20}{km}$};
\end{tikzpicture}

{\small Le voyage du muon~: une durée de vie classique s'éteint à
\qty{660}{m} sous l'altitude de naissance~; la durée de vie dilatée
traverse l'atmosphère.}
\end{omfigure}

Les confirmations s'accumulent. En 1971 Hafele et Keating ont fait
voler des horloges atomiques au césium (\cref{ch:g12:oscillators-and-time})
autour du monde sur des avions de ligne~; à l'atterrissage elles
différaient des horloges au sol des fractions de microseconde
prédites. Les accélérateurs misent sur la dilatation chaque jour,
leurs particules poussées à des $\gamma$ de milliers survivant des
tours qu'elles ne pourraient jamais classiquement achever. Et le GPS
est une expérience en cours avec une étiquette de prix.

\begin{remark}[La facture relativiste du GPS]\label{rem:g12:special-relativity:gps}
Un \omterm{def:g10:universal-gravitation:satellite}{satellite} GPS
\omterm{def:g10:universal-gravitation:satellite}{orbite} à environ \qty{3.9}{km/s}, donc la
\omterm{thm:g12:special-relativity:postulates}{relativité restreinte} ralentit son horloge de
quelque \qty{7}{\micro s} par jour. La gravité, plus faible en
altitude, \omterm{def:g11:work-of-force:work}{travaille} dans l'autre sens et gagne~: la
relativité générale (la théorie d'Einstein de la
\omterm{def:g11:fundamental-interactions:four}{gravitation}, racontée dans les volumes
universitaires) l'accélère d'environ \qty{45}{\micro s} par jour~; net,
\qty{38}{\micro s} par jour d'avance. Le positionnement convertit les
lectures d'horloge en distances à la vitesse $c$~: une journée non
corrigée étalerait les positions de $c \times \qty{38}{\micro s}
\approx \qty{11}{km}$ --- d'où les horloges désaccordées du paragraphe
d'ouverture.
\end{remark}

\begin{proposition}[Contraction des longueurs]\label{prop:g12:special-relativity:contraction}
\index{contraction des longueurs}
Un corps de longueur $L_0$ dans son référentiel de repos est mesuré
plus court le long de sa direction de mouvement dans un référentiel
où il se déplace à la vitesse $v$~: $L = L_0/\gamma$.
\end{proposition}
\admitted

\begin{remark}[L'histoire propre du muon]\label{rem:g12:special-relativity:muonframe}
La \omterm{prop:g12:special-relativity:contraction}{contraction des longueurs} est le revers de
la \omterm{thm:g12:special-relativity:dilation}{dilatation du temps} (dérivée honnêtement dans
les volumes universitaires)~: l'horloge du muon tic-tac normalement
dans son propre référentiel, mais l'atmosphère qui monte à sa
rencontre est contractée à $\qty{15}{km}/30 = \qty{500}{m}$, traversée
bien dans \qty{2.2}{\micro s} --- les deux référentiels s'accordent
sur le fait que le muon atteint le sol.
\end{remark}

\section{Masse, énergie, et le monde classique retrouvé}

\begin{theorem}[Équivalence masse--énergie]\label{thm:g12:special-relativity:emc2}
\index{équivalence masse--énergie}
Un corps de masse $m$ possède, par sa masse seule,
l'\omterm{def:g10:energy-conservation:energy}{énergie}
\[ E = m c^2 . \]
La masse est une forme concentrée d'\omterm{def:g10:energy-conservation:energy}{énergie}~;
l'\omterm{def:g10:energy-conservation:energy}{énergie}, inversement, a de la masse.
\end{theorem}
\admitted

\begin{remark}[Où vous l'avez rencontrée]\label{rem:g12:special-relativity:emc2seen}
C'est le grand livre derrière l'\omterm{def:g10:energy-conservation:forms}{énergie nucléaire}
(\cref{ch:g12:nuclear-energy})~: le \omterm{def:g12:nuclear-energy:defect}{défaut de masse}
d'une réaction, multiplié par $c^2$, est l'\omterm{def:g10:energy-conservation:energy}{énergie}
libérée. Un gramme de masse vaut \qty{9.0e13}{J} --- une journée
entière de production d'une grande centrale de
\omterm{def:g11:work-of-force:power}{puissance}. La dérivation honnête appartient aux
volumes universitaires~; les réacteurs et les étoiles en vivent
quoi qu'il en soit.
\end{remark}

\begin{remark}[Newton n'avait pas tort]\label{rem:g12:special-relativity:newton}
Poser $v \ll c$~: $\gamma \to 1$, les temps s'accordent, les longueurs
ne se contractent pas, la simultanéité redevient absolue --- la
mécanique des trois dernières années réapparaît, intacte, comme la
limite bas vitesse. La relativité n'a pas démoli Newton~; elle a
tracé la frontière de son empire. C'est ainsi que la physique
grandit~: en affinant les théories, chacune rendant l'ancienne comme
cas limite.
\end{remark}

\section{Exercices}

\begin{exercise}[$\star$]\label{exo:g12:special-relativity:1}
Calculer $\gamma$, à trois
\omterm{def:g10:orders-of-magnitude:sigfig}{chiffres significatifs}, pour $v = 0.1c$,
$0.6c$, $0.8c$, $0.995c$.
\end{exercise}

\begin{exercise}[$\star$]\label{exo:g12:special-relativity:2}
Un vaisseau croise à $0.6c$~; son métronome bat toutes les
\qty{10.0}{s} selon l'horloge du vaisseau. Quel intervalle le sol
mesure-t-il~? Lequel est le \omterm{def:g12:special-relativity:proper-time}{temps propre}, et
pourquoi~?
\end{exercise}

\begin{exercise}[$\star$]\label{exo:g12:special-relativity:3}
À quelle vitesse (en fraction de $c$) a-t-on $\gamma = 2$~? Et
$\gamma = 10$~?
\end{exercise}

\begin{exercise}[$\star$]\label{exo:g12:special-relativity:4}
Une voiture roule à \qty{130}{km/h}. Calculer $v/c$, puis
$\gamma - 1$ (\cref{rem:g12:special-relativity:everyday})~; combien de
temps doit-elle rouler pour que son horloge retarde de \qty{1}{s}~?
\end{exercise}

\begin{exercise}[$\star$]\label{exo:g12:special-relativity:5}
Dans l'expérience du train d'Einstein, expliquer --- en citant le
postulat utilisé --- pourquoi le passager conclut que la frappe
\emph{avant} est venue la première, et pourquoi aucun observateur
n'a tort.
\end{exercise}

\begin{exercise}[$\star\star$]\label{exo:g12:special-relativity:6}
Des pions chargés vivent $\Delta t_0 = \qty{2.6e-8}{s}$. Un faisceau
quitte un accélérateur à $0.99c$. Calculer $\gamma$, la durée de vie
dans le labo, la \omterm{def:g12:projectile-motion:range}{portée} classique $v\,\Delta t_0$,
et la \omterm{def:g12:projectile-motion:range}{portée} moyenne réelle.
\end{exercise}

\begin{exercise}[$\star\star$]\label{exo:g12:special-relativity:7}
Sur la courbe de $\gamma$, lire les vitesses donnant $\gamma = 1.5$ et
$\gamma = 3$ (vérifier par le calcul). Que dit l'asymptote verticale
en $v = c$ sur l'accélération d'un corps massif~?
\end{exercise}

\begin{exercise}[$\star\star$]\label{exo:g12:special-relativity:8}
Une horloge à lumière a des miroirs distants de $L = \qty{1.5}{m}$.
Calculer son tic au repos, puis son tic mesuré depuis le sol lorsqu'elle
croise à $0.8c$.
\end{exercise}

\begin{exercise}[$\star\star$]\label{exo:g12:special-relativity:9}
Qui mesure le \omterm{def:g12:special-relativity:proper-time}{temps propre} entre~: (a)~deux
tics d'une montre --- le porteur, ou un passant~? (b)~naissance et
désintégration d'un muon --- le muon, ou le labo~? (c)~quitter la
Terre et atteindre Mars --- le pilote, ou le centre de mission~?
\end{exercise}

\begin{exercise}[$\star\star$]\label{exo:g12:special-relativity:10}
La station spatiale \omterm{def:g10:universal-gravitation:satellite}{orbite} à \qty{7.7}{km/s}.
Calculer $\gamma - 1 \approx v^2/2c^2$, puis de combien l'horloge d'un
astronaute retarde sur une horloge au sol après six mois
(\qty{1.58e7}{s}).
\end{exercise}

\begin{exercise}[$\star\star$]\label{exo:g12:special-relativity:11}
Calculer le contenu d'\omterm{def:g10:energy-conservation:energy}{énergie} $E = mc^2$ de
\qty{1.0}{g} de matière. Le comparer à la production quotidienne d'une
centrale de \omterm{def:g11:work-of-force:power}{puissance} de \qty{1.0}{GW}, et à
l'\omterm{def:g11:mechanical-energy:kinetic}{énergie cinétique} de \qty{8.5e5}{J} d'une
voiture à \qty{130}{km/h}.
\end{exercise}

\begin{exercise}[$\star\star\star$]\label{exo:g12:special-relativity:12}
Montrer à partir de $\gamma = 1/\sqrt{1 - v^2/c^2}$ que
$v = c\sqrt{1 - 1/\gamma^2}$~; en déduire la vitesse des muons à
$\gamma = 30$, à cinq \omterm{def:g10:orders-of-magnitude:sigfig}{chiffres significatifs}
en \omterm{def:g10:orders-of-magnitude:unit}{unités} de $c$.
\end{exercise}

\begin{exercise}[$\star\star\star$]\label{exo:g12:special-relativity:13}
Raconter le voyage du muon depuis son propre référentiel~: quelle
épaisseur a l'atmosphère à $\gamma = 30$, combien de temps la
traverse-t-il à $0.99944c$, et cela tient-il dans une durée de vie
propre~?
\end{exercise}

\begin{exercise}[$\star\star\star$]\label{exo:g12:special-relativity:14}
Un accélérateur linéaire pousse des électrons à $\gamma = \num{1.0e5}$
le long d'un tube de \qty{3.2}{km}. Combien dure le trajet dans le
labo~? Combien dans le référentiel de l'électron, et quelle
«~longueur~» a le tube là~?
\end{exercise}

\begin{exercise}[$\star\star\star$]\label{exo:g12:special-relativity:15}
Proxima du Centaure se trouve à $4.2$
\omterm{def:g10:orders-of-magnitude:lightyear}{années-lumière} (une
\omterm{def:g10:orders-of-magnitude:lightyear}{année-lumière} est la distance que la
lumière parcourt en un an). Une sonde y croise à $0.9c$~: combien
dure le trajet sur le calendrier de la Terre, et sur l'horloge de la
sonde~?
\end{exercise}

\section{Problème~: Le muon et le jumeau}

\begin{problem}[{Problème du week-end --- le voyage du muon et le
jumeau de l'astronaute~: une particule qui survit à son horloge, une
horloge à lumière reconstruite, un voyageur plus jeune que son jumeau,
et la comptabilité qui tient le GPS honnête}]
\label{pb:g12:special-relativity:1}
Données~: durée de vie propre du muon $\Delta t_0 = \qty{2.2}{\micro s}$~;
muons nés à environ \qty{15}{km} de hauteur à près de $c$~; flux au
niveau de la mer d'environ un muon par \unit{cm^2} par minute~;
$c = \qty{3.00e8}{m/s}$.

\medskip
\noindent\textbf{Partie I --- L'arrivée impossible du muon.}
\begin{enumerate}
  \item Sans relativité, quelle distance un muon parcourt-il en une
        durée de vie~?
  \item Combien de durées de vie la descente de \qty{15}{km} prend-elle
        alors~?
  \item Après $n$~durées de vie une fraction $\mathrm{e}^{-n}$
        survit~: l'estimer, confronter le flux mesuré, et passer
        verdict sur la physique classique en une phrase.
  \item Ces muons arrivent en fait avec $\gamma \approx 30$. Calculer
        leur durée de vie mesurée depuis le sol et la
        \omterm{def:g12:projectile-motion:range}{portée} qui en résulte.
  \item Calculer leur vitesse, en utilisant
        $v = c\sqrt{1 - 1/\gamma^2}$, à cinq
        \omterm{def:g10:orders-of-magnitude:sigfig}{chiffres significatifs} en
        \omterm{def:g10:orders-of-magnitude:unit}{unités} de $c$.
\end{enumerate}

\medskip
\noindent\textbf{Partie II --- L'horloge à lumière, reconstruite.}
Une horloge à lumière ($L = \qty{1.5}{m}$) voyage sur un vaisseau à
$v = 0.6c$.
\begin{enumerate}[resume]
  \item Calculer le tic $\Delta t_0$ lu à bord.
  \item Vu du sol, un tic dure $\Delta t$~: donner l'avance horizontale
        par demi-tic, la montée verticale du
        \omterm{def:g12:quantum-world:photon}{photon}, et --- en nommant le postulat qui la
        fixe --- la longueur du demi-chemin oblique.
  \item Appliquer Pythagore au triangle du demi-tic et résoudre pour
        $\Delta t$, en montrant que $\Delta t = \gamma\,\Delta t_0$.
  \item Calculer $\gamma$ et $\Delta t$ numériquement.
  \item Quel observateur mesure le
        \omterm{def:g12:special-relativity:proper-time}{temps propre} entre deux tics~?
        Justifier à partir de la définition.
\end{enumerate}

\medskip
\noindent\textbf{Partie III --- Le voyageur et son jumeau.}
Une astronaute quitte son jumeau sur Terre pour un aller-retour à
$0.8c$~; les horloges terrestres chronométrent tout le voyage à
\qty{10.0}{years}.
\begin{enumerate}[resume]
  \item Calculer $\gamma$ à $0.8c$ (fraction exacte, puis décimal).
  \item Combien de temps s'écoule sur l'horloge de l'astronaute~?
  \item Quelle distance (référentiel Terre, en
        \omterm{def:g10:orders-of-magnitude:lightyear}{années-lumière}) le demi-tour a-t-il
        atteint~?
  \item Quelle est la différence d'âge aux retrouvailles, et quel
        jumeau est le plus jeune~?
  \item «~Mais le mouvement est relatif --- l'astronaute pourrait
        prétendre que la Terre a voyagé~!~» Résoudre l'objection en
        une phrase.
\end{enumerate}

\medskip
\noindent\textbf{Partie IV --- La facture GPS.}
Un \omterm{def:g10:universal-gravitation:satellite}{satellite} GPS
\omterm{def:g10:universal-gravitation:satellite}{orbite} au rayon $r = \qty{2.66e7}{m}$,
\omterm{def:g12:oscillators-and-time:oscillator}{période} $T = \qty{12}{h}$.
\begin{enumerate}[resume]
  \item Calculer sa vitesse orbitale $v = 2\pi r/T$.
  \item Calculer $\gamma - 1 \approx v^2/2c^2$.
  \item En déduire, en microsecondes, de combien l'horloge du
        \omterm{def:g10:universal-gravitation:satellite}{satellite} retarde sur les horloges
        au sol par jour (\qty{86400}{s}).
  \item Si cette seule dérive n'était pas corrigée, les positions
        errent de $c \times$ (dérive)~: calculer l'erreur après un
        jour.
  \item La gravité, plus faible en altitude, fait tourner l'horloge
        \emph{vite} de \qty{45}{\micro s} par jour (relativité
        générale)~: calculer la dérive nette et l'erreur de position
        quotidienne en kilomètres --- le nombre qui oblige les
        ingénieurs à désaccorder chaque horloge GPS.
\end{enumerate}
\end{problem}

\bigskip
\noindent
Ici s'achève la physique de l'école --- pas la physique. Dans les
volumes universitaires, la mécanique revient armée du calcul
différentiel~; l'électricité et le magnétisme se révèlent comme une
seule structure dans laquelle la lumière, et la relativité
elle-même, se cachait depuis le début~; le monde quantique, entrevu
au chapitre précédent, devient une théorie opérationnelle de la
matière. Chaque théorie sera tenue, comme la relativité, de rendre
ce que vous savez maintenant comme cas limite. Emportez-le~: c'est
la part qui ne changera pas.
